\section{Appendix W: TRXT Nullivance Parameter Dictionary (Audit V5.4)}
\label{app:dictionary}

The model is governed by the following strict set of parameters.  We distinguish clearly between \textbf{inputs} (experimentally measured constants), \textbf{derived} quantities (determined by the theoretical chain with no free tuning), and \textbf{fit} parameters (adjusted to match observations).

\subsection{Parameter Classification and Derivation Status}

\begin{table}[H]
\centering
\small
\begin{tabular}{llccl}
\toprule
\textbf{Symbol} & \textbf{Physical Meaning} & \textbf{Value} & \textbf{Status} & \textbf{Derivation Source} \\
\midrule
\multicolumn{5}{l}{\textit{Tier 0: Fundamental inputs (measured, not derived)}} \\
$G$ & Newton Constant & $6.674 \times 10^{-11}$ & Input & Cavendish experiment \\
$\alpha$ & Fine Structure Constant & $1/137.036$ & Input & QED measurement \\
$m_\tau$ & Tau Lepton Mass & $1776.86$ MeV & Input* & PDG 2024 \\
\midrule
\multicolumn{5}{l}{\textit{Tier 1: Derived from Tier 0 (no free parameters)}} \\
$M_{\mathrm{Pl}}$ & Planck Mass & $1.22 \times 10^{19}$ GeV & Derived & $\sqrt{\hbar c / G}$ \\
$M^*$ & Condensate Scale & $365.24 \pm 0.02$ GeV & \textbf{Derived} & $m_\tau \times 3/(2\alpha)$ \\
$\xi$ & Non-minimal Coupling & $\approx 1/6$ & Derived & Conformal invariance \\
\midrule
\multicolumn{5}{l}{\textit{Tier 2: Fit or constrained parameters}} \\
$n$ & Polytropic Index & $1.37 \pm 0.1$ & Best Fit & DM halo profile matching \\
$\Lambda$ & Cosmological Constant & $10^{-52} \text{ m}^{-2}$ & Constrained & NJL condensate (see C4) \\
$c_s$ & Sound Speed (Core) & $< 10^{-3}$ & Constraint & EOS stability \\
$\omega^2$ & Stability Eigenvalue & $\ge 0$ & Verified & Ghost-free check \\
\bottomrule
\end{tabular}
\caption{Unified Parameter Dictionary V5.4.  *See Section~W.2 for the circularity analysis of $m_\tau$ vs $M^*$.}
\label{tab:params}
\end{table}

\subsection{W.1: Error Budget Analysis}
\label{sec:error_budget}
We propagate the experimental uncertainties of input parameters to derive the precision of $M^*$.

\begin{table}[H]
\centering
\begin{tabular}{l c c}
\toprule
\textbf{Input Parameter} & \textbf{Value (PDG 2024)} & \textbf{Relative Error} \\
\midrule
$\alpha^{-1}$ & $137.035999084(21)$ & $1.5 \times 10^{-10}$ \\
$m_\tau$ & $1776.86 \pm 0.12$ MeV & $6.7 \times 10^{-5}$ \\
\midrule
\textbf{Derived Quantity} & \textbf{Result} & \textbf{Propagated Error} \\
\midrule
$M^*$ & $365.24$ GeV & $\pm 25$ MeV \\
\bottomrule
\end{tabular}
\caption{Error Budget. The precision is dominated by the tau mass measurement.}
\end{table}

\subsection{W.2: Circularity Analysis --- $M^*$ vs $m_\tau$}
\label{sec:circularity}

\textbf{Honest assessment of the input/output chain:}

The relation $M^* = m_\tau \times 3/(2\alpha)$ uses the measured tau mass as input.  The paper's BCS gap equation chain:
\begin{equation}
    \alpha \;\to\; X = \frac{3}{2\alpha} \;\to\; q = 6 \text{ (Abrikosov)} \;\to\; k_F = \frac{5}{6} \;\to\; \eta \;\to\; t \;\to\; C = 5.339 \;\to\; M^*
\end{equation}
requires the Fermi velocity $v_F$ to close; the value $v_F = 1/5$ reproduces $M^* = 365.24$~GeV but is not yet derived from first principles.

\textbf{What IS derivable:}
\begin{itemize}
    \item $k_F = 5/6$ from Abrikosov~\cite{abrikosov1957} vortex lattice energy minimization ($C_6$ symmetry, $q=6$). \textbf{Not a fit.}
    \item $g = 4$ from Dirac band structure degeneracy (spin $\times$ valley). \textbf{Not a fit.}
    \item $N_f = 16$ from $\mathrm{Cl}(6)$ minimal left ideal dimension. \textbf{Not a fit.}
    \item $\Lambda_{\mathrm{UV}}$ from induced gravity heat kernel. \textbf{Not a fit.}
\end{itemize}

\textbf{What is NOT yet derivable:}
\begin{itemize}
    \item $v_F = 1/5$: the Fermi velocity in the tight-binding model.  The BCS self-consistency equation determines $v_F$ \emph{given} $M^*$, but the reverse direction (predicting $M^*$ from $v_F$) requires a microscopic lattice computation not yet performed.
\end{itemize}

\textbf{Current honest status:} The TRXT framework is a \textbf{one-parameter theory} with input $m_\tau$ (equivalently $M^*$). Given this single input, it predicts the W, Z, and Higgs masses to $< 0.2\%$ accuracy.  The BCS chain provides a plausible \emph{route} to eliminating this last input, contingent on deriving $v_F$ from the underlying lattice model.  Until $v_F$ is independently derived, $m_\tau$ should be classified as an \textbf{input}, not ``Derived (V1)''.
